%! Advanced Factoring — Unit 3, Lesson 3.2 — Group Activity
\documentclass[10pt]{article}
\usepackage{algebra2-article}
\usepackage{algebra2-boxes}
\newcommand{\TallMath}[1]{$\displaystyle #1\rule[-1.4em]{0pt}{3.2em}$}

\begin{document}

\pageheader{Unit 3, Lesson 3.2}{Group Activity}
\namepartnerperiod

\vspace{0.08in}

\begin{headlinebox}{forestbg}
\small\textbf{Taking Polynomials Apart.} Yesterday you built polynomials by multiplying; today you
take them apart. Two habits: \emph{GCF first, every time}, then \emph{count the terms} to pick your
tool. And one standard --- \textbf{completely}. Before you call a problem finished, look at every
factor and ask ``can this still be broken apart?'' You can always check by multiplying back.
\end{headlinebox}

\vspace{0.06in}

\begin{tcolorbox}[colback=white, colframe=black!40, title=\textbf{Tier R --- Remediate},
    fonttitle=\bfseries, arc=1.5mm, left=3mm, right=3mm, top=2mm, bottom=2mm, breakable]
\begin{enumerate}[leftmargin=1.7em, itemsep=9pt]
  \item \textbf{GCF.} $8x^4 - 20x^3 + 14x^2 = $ \blank{5.5cm}

    \vspace{2pt}
    Check: multiply your answer back out. Did you get the original? \blank{1.4cm}

  \item \textbf{GCF, two variables.} $9x^2y - 15xy^2 = $ \blank{5.5cm}

  \item \textbf{Two terms.} $x^2 - 36 = $ \blank{5.0cm}

    \vspace{2pt}
    Which pattern was that? \blank{4.4cm}

  \item \textbf{Three terms.} $x^2 + 9x + 20 = $ \blank{5.0cm}

  \item \textbf{Four terms --- grouping, step by step.} $x^3 + 2x^2 + 5x + 10$

    \vspace{3pt}
    \emph{Step 1 --- the two pairs:} \; $(x^3 + 2x^2) \; + \; (5x + 10)$

    \vspace{3pt}
    \emph{Step 2 --- GCF each pair:} \; \blank{2.0cm}$(x+2)$ \; $+$ \; \blank{2.0cm}$(x+2)$

    \vspace{3pt}
    \emph{Step 3 --- factor out the shared binomial:} \blank{4.6cm}

    \vspace{3pt}
    \emph{Step 4 --- is $x^2+5$ factorable?} \blank{1.4cm} \; So the final answer is
    \blank{4.0cm}
\end{enumerate}
\end{tcolorbox}

\begin{tcolorbox}[colback=white, colframe=black!40, title=\textbf{Tier A --- Approaching Proficiency},
    fonttitle=\bfseries, arc=1.5mm, left=3mm, right=3mm, top=2mm, bottom=2mm, breakable]
\begin{enumerate}[leftmargin=1.7em, itemsep=10pt]
  \item \textbf{Two steps --- GCF, then a pattern.}

    \vspace{2pt}
    (a) $3x^3 - 75x = $ \blank{5.0cm} \qquad (b) $4x^4 - 64 = $ \blank{5.0cm}

    \vspace{2pt}
    In (b), how many factors does your final answer have? \blank{1.0cm} \; Is one of them prime?
    Which, and why? \writeline

  \item \textbf{Perfect square, two variables.} $9x^2 - 30xy + 25y^2 = $ \blank{4.6cm}

    \vspace{2pt}
    Verify the middle term: $2 \cdot$ \blank{1.2cm} $\cdot$ \blank{1.2cm} $= $ \blank{1.6cm}

  \item \textbf{Cubes.} Name $a$ and $b$ \emph{before} you write anything else.

    \vspace{3pt}
    (a) $x^3 + 27$: \quad $a = $ \blank{0.9cm}, \; $b = $ \blank{0.9cm} \quad $\Rightarrow$
    \blank{4.6cm}

    \vspace{3pt}
    (b) $64x^3 - 1$: \quad $a = $ \blank{0.9cm}, \; $b = $ \blank{0.9cm} \quad $\Rightarrow$
    \blank{4.6cm}

  \item \textbf{Three terms, $a \ne 1$ --- by grouping.} $3x^2 - 10x - 8$

    \vspace{2pt}
    Two numbers with product $3 \cdot (-8) = -24$ and sum $-10$: \blank{1.2cm} and \blank{1.2cm}

    \vspace{2pt}
    Rewrite the middle term and group: \blank{6.0cm}

    \vspace{2pt}
    Answer: \blank{4.4cm}

  \item \textbf{Prime or not?} Decide, and justify each in one line.

    \vspace{2pt}
    (a) $x^2 + 9$ \quad \blank{1.6cm} \; because \blank{6.0cm}

    \vspace{2pt}
    (b) $x^3 + 27$ \quad \blank{1.6cm} \; because \blank{6.0cm}

  \item \textbf{Error analysis.} A classmate turns in $x^4 - 16 = (x^2 - 4)(x^2 + 4)$ as a final
    answer and is marked wrong.

    \vspace{2pt}
    (a) Is anything they wrote \emph{incorrect}? \blank{1.4cm}

    \vspace{2pt}
    (b) So why is it wrong? Write the complete factorization: \blank{5.0cm}

    \vspace{2pt}
    (c) In one sentence, tell them what to check before turning in a factoring problem. \writeline
\end{enumerate}
\end{tcolorbox}

\begin{tcolorbox}[colback=white, colframe=black!40, title=\textbf{Tier E --- Extension},
    fonttitle=\bfseries, arc=1.5mm, left=3mm, right=3mm, top=2mm, bottom=2mm, breakable]
\textbf{Part 1 --- Closing the loop on the graph.} In Lesson 3.0 you read the graph of
$q(x) = x^4 - 5x^2 + 4$ and found four $x$-intercepts.
\begin{enumerate}[label=(\alph*), leftmargin=1.8em, itemsep=6pt]
  \item It has three terms in \textbf{quadratic form}. Factor it as though $x^2$ were the variable:
        \blank{5.0cm}
  \item Finish the job: \; $q(x) = $ \blank{6.0cm}
  \item List the zeros: \blank{4.0cm} \; Each is a zero of multiplicity \blank{1.0cm}, so the graph
        \blank{2.6cm} the $x$-axis at every one of them --- which is exactly what you saw.
  \item The degree is $4$ and the leading coefficient is $1$. Multiply the four leading terms of your
        factors together to confirm the leading term is $x^4$: \blank{3.4cm}
\end{enumerate}

\vspace{5pt}
\textbf{Part 2 --- Squares before cubes.} The binomial $x^6 - 64$ is \emph{both} a difference of
squares ($x^6 = (x^3)^2$, $64 = 8^2$) and a difference of cubes ($x^6 = (x^2)^3$, $64 = 4^3$). Try
both roads.
\begin{enumerate}[label=(\alph*), leftmargin=1.8em, itemsep=6pt]
  \item \textbf{Road 1 --- squares first.} \; $x^6 - 64 = ($\blank{1.6cm}$)($\blank{1.6cm}$)$

        \vspace{2pt}
        Now factor each of those as cubes: \blank{6.5cm}

  \item \textbf{Road 2 --- cubes first.} \; $x^6 - 64 = ($\blank{1.6cm}$)($\blank{2.6cm}$)$

        \vspace{2pt}
        Factor the first of those two: \blank{3.6cm} \; Can you factor the second one with today's
        tools? \blank{1.4cm}

  \item Road 2 stalls on $x^4 + 4x^2 + 16$. (For the record, it equals $(x^2+2x+4)(x^2-2x+4)$ ---
        multiply those two together to check.) \blank{5.0cm}

  \item Both roads reach the same four factors. Which road got there without help, and what is the
        rule of thumb? \writelines{2}
\end{enumerate}

\vspace{5pt}
\textbf{Part 3 --- Prove the identity, then hit the wall.}
\begin{enumerate}[label=(\alph*), leftmargin=1.8em, itemsep=6pt]
  \item Multiply $(a-b)(a^2+ab+b^2)$ out in full --- all six products --- and show that four of them
        cancel, leaving $a^3 - b^3$. \writelines{3}
  \item Now try today's entire toolkit on the unit's anchor, $x^3 - 3x + 2$. GCF? \blank{1.2cm} \;
        A pattern for two terms? \blank{1.2cm} \; Quadratic form? \blank{1.2cm} \; Four terms to
        group? \blank{1.2cm}
  \item You already know from Lesson 3.0 that it equals $(x-1)^2(x+2)$. So what does that tell you
        about the toolkit --- and what would you need in order to find that factorization yourself?
        \writelines{2}
\end{enumerate}
\end{tcolorbox}

\end{document}
